\documentclass[10pt,a4paper]{amsart}
\usepackage[margin=19mm]{geometry}
\usepackage{amsmath,amssymb,mathtools}
\usepackage[scr=rsfs,cal=euler]{mathalfa}
\usepackage{xfrac}
\usepackage[unicode,hidelinks,bookmarks=false]{hyperref}
\newcommand{\ie}{i.e.\ }
\newcommand{\cf}{cf.\ }
\newcommand{\resp}{respectively\ }
\newcommand{\Subsection}{\subsection}
\providecommand{\nobreakdash}{\nobreak\hbox{-}}
\setlength{\emergencystretch}{2em}
\multlinegap0pt
\usepackage{amssymb}
\usepackage{mathtools}
\newtheorem{theorem}{Theorem}
\title{Arbitrarily Fast Quantum Dispersion \\ in Long-Range Crystals}
\author{Ga\'etan Leclerc, Mostafa Sabri, Tuomas Sahlsten}
\date{Source: arXiv:2608.27326v1 (CC BY 4.0)}
\begin{document}
\maketitle

\noindent\emph{Source and licence.} Arbitrarily Fast Quantum Dispersion \\ in Long-Range Crystals. Version arXiv:2608.27326v1 (CC BY 4.0), distributed by arXiv under the Creative Commons Attribution 4.0 International licence, http://creativecommons.org/licenses/by/4.0/, as recorded in the arXiv licence metadata for this exact version and shown on the abstract page https://arxiv.org/abs/2608.27326v1. Full licence notice: \texttt{licenses/arxiv-2608.27326-CC-BY-4.0.txt}.\par\smallskip

\noindent\textit{Editorial scope.} Counted unit: the article's theorem theorem (source0.tex lines 627--635). The article's complete original proof of it is delivered with it (source0.tex lines 638--678, one \texttt{proof} environment of 41 lines). No mathematics is written, repaired or re-derived: the statement and the proof are the article's own lines, in the article's own notation, and no retained line is substituted or re-flowed.  No further source-local context is needed: every cross-reference of every retained line resolves inside the two delivered spans.  Nothing was omitted from any retained span, so the package carries no omission record.  No retained line contains a citation, so the wrapper carries no bibliography and no bibliography entry.  No retained line contains a figure environment, an image-inclusion command or drawing code, so no image is delivered and the package declares no asset dependency.  Editorial formatting only, in addition: an AMS \texttt{amsart} preamble that restates the article's own declarations and macros used by the retained lines, namely the environments \texttt{theorem} on separate counters, together with the standard packages those lines need.  The retained environments print in this document as Theorem 1 in order of appearance, whereas the article prints the same items with the numbering of its own full document.  The complete native source of the article as distributed in the arXiv e-print for this exact version is bundled unchanged as \texttt{sources/208708/source0.tex}.
\begin{theorem}
There exists $\rho \in (0,1)$ that depends only on $\varphi$ such that, for all $\beta \in (0,1/2]$, for all $\lambda$ sufficiently large and all $t$ sufficiently large depending on $\beta,\lambda$,, we have, uniformly in $m$,
$$
\forall h \in C^1, \qquad
\int_{\mathbb{S}^1} |\mathcal{L}_{\lambda,i\Phi(t,m)}^{\langle n_\beta(t) \rangle} h|^2 d\theta
\leq \frac{\|h\|_{1,t}^2}{t^{\rho}}.
$$
The condition on $\lambda$ is of the form $\lambda \geq \max(C_1(\varphi),\mu^{-C_2(\varphi)/\beta}).$
\end{theorem}

\begin{proof}[Proof (of $L^2$-Dolgopyat estimates)]
Let $\beta \in (0,1/2]$ and recall that $n := n_\beta(t)$ is such that $\mu^{n_\beta(t)} \sim t^{-\beta} \geq t^{-1/2}$. We suppose that $\lambda$ is large enough so that the previous tree lemma applies, and that $\lambda \geq \mu^{-10/\beta}.$ We start by writing (denoting $\Phi = \Phi(t,m)$ for clarity):
\begin{align*} 
\Big\| \mathcal{L}_{\lambda,\Phi(t,m)}^{\langle n \rangle}(h) \Big\|_{L^2([0,1])}^2 &= \int_0^1 |\mathcal{L}_{\lambda,i\Phi(t,m)}^{\langle n \rangle}(h)(\theta)|^2 d\theta \\
& = \int_0^1 \Big| \lambda^{-n} \sum_{\mathbf{a}} e^{i S_n \Phi \circ g_{\mathbf{a}}(\theta)} h(g_\mathbf{a})(\theta) \Big|^2 d\theta \\
& = \lambda^{-2n} \sum_{\mathbf{a},\mathbf{b} \in \mathcal{A}^n(\lambda)} \int_0^1 e^{i \big( S_n \Phi \circ g_\mathbf{a}(\theta) - S_n \Phi \circ g_\mathbf{b}(\theta) \big)} h(g_\mathbf{a}(\theta)) \overline{h(g_\mathbf{b}(\theta))} d\theta.
\end{align*}
Let us then iterate a little more our transfer operator: recall that we have the relation $\int h d\theta = \int \mathcal{L}_\lambda^{\mathfrak{n}}(h) d\theta$ for any $\mathfrak{n}$. We then choose $\mathfrak{n}$ such that $\lambda^{2\mathfrak{n}} \simeq t^{1/8}$ and we write:
$$ \Big\| \mathcal{L}_{\lambda,\Phi(t,m)}^{\langle n \rangle}(h) \Big\|_{L^2([0,1])}^2 = \frac{1}{\lambda^{2n+2 \mathfrak{n}}} \underset{\mathbf{c} \in \mathcal{A}(\lambda)^{2\mathfrak{n}}}{\sum_{\mathbf{a},\mathbf{b} \in \mathcal{A}(\lambda)^n}} \int_0^1 e^{i \big( S_n \Phi \circ g_\mathbf{a}(g_\mathbf{c}(\theta)) - S_n \Phi \circ g_\mathbf{b}(g_\mathbf{c}(\theta)) \big)} h(g_\mathbf{ac}(\theta)) \overline{h(g_\mathbf{bc}(\theta))} d\theta. $$
Since $h$ is $C^1$ and $g_{\mathbf{ac}}' = \lambda^{-n-2\mathfrak{n}} \leq t^{-2}$ for $t$ large enough (as $\lambda \geq \mu^{-10/\beta}$),, this gives, writing $h(g_{\mathbf{ac}}(\theta)) = h(x_{\mathbf{ac}})+\mathcal{O}(\|h\|_{C^1} t^{-2})$:
$$ \Big\| \mathcal{L}_{\lambda,\Phi(t,m)}^{\langle n \rangle}(h) \Big\|_{L^2([0,1])}^2 \leq \frac{\|h\|_\infty^2}{\lambda^{2n+2\mathfrak{n}}} \underset{\mathbf{c} \in \mathcal{A}(\lambda)^{2\mathfrak{n}}}{\sum_{\mathbf{a},\mathbf{b} \in \mathcal{A}(\lambda)^n}} \Big| \int_0^1 e^{i \big( S_n \Phi \circ g_\mathbf{a}(g_\mathbf{c}(\theta)) - S_n \Phi \circ g_\mathbf{b}(g_\mathbf{c}(\theta)) \big)}  d\theta \Big| + \frac{\|h\|_{{1,t}}^2}{t} .$$
Let us then concentrate on the integral $I_{\mathbf{a},\mathbf{b},\mathbf{c}}(t) := \int_0^1 e^{i \big( S_n \Phi \circ g_\mathbf{a}(g_\mathbf{c}(\theta)) - S_n \Phi \circ g_\mathbf{b}(g_\mathbf{c}(\theta)) \big)}  d\theta$. We have
$$ I_{\mathbf{a},\mathbf{b},\mathbf{c}}(t) =  \int_0^1 e^{i T \psi_{\mathbf{a},\mathbf{b},\mathbf{c}}(\theta)}  d\theta $$
with $T=t \mu^{n-1} \lambda^{-5\mathfrak{n}-2} \geq t^{1/8}$ and 
$$\psi_{\mathbf{a},\mathbf{b},\mathbf{c}}(\theta) = t^{-1} \mu^{1-n} \lambda^{5\mathfrak{n}+2} \big( S_n \Phi \circ g_\mathbf{a}(g_\mathbf{c}(\theta)) - S_n \Phi \circ g_\mathbf{b}(g_\mathbf{c}(\theta)) \big).$$

Again the classical van Der Corput lemma ensures that $|I_{\mathbf{a},\mathbf{b},\mathbf{c}}(t)| \leq c T^{-1/2}$ (for some universal constant $c$) if $|\psi_{\mathbf{a},\mathbf{b},\mathbf{c}}^{(2)}(\theta)| \geq 1$ on $[0,1]$. We hence get the bound:
$$ \Big\| \mathcal{L}_{\lambda,\Phi(t,m)}^{\langle n \rangle}(h) \Big\|_{L^2([0,1])}^2 \lesssim \|h\|_{1,t}^2 \Big( \frac{1}{t^{1/16}} + \lambda^{-2n-2\mathfrak{n}} \#\{ (\mathbf{a},\mathbf{b},\mathbf{c}) \in \mathcal{A}^{2n+2\mathfrak{n}} \ | \ \exists \theta, |\psi_{\mathbf{a},\mathbf{b},\mathbf{c}}^{(2)}(\theta)| \leq 1\} \Big). $$
To conclude, we need to bound this cardinality. We first compute:
\begin{align*}
\psi_{\mathbf{a},\mathbf{b},\mathbf{c}}^{(2)}(\theta) &=: \lambda^{\mathfrak{n}} \big( \ell_{\mathbf{a}}(g_\mathbf{c}(\theta)) - \ell_{\mathbf{b}}(g_\mathbf{c}(\theta)) \big)  \\
& = \lambda^{\mathfrak{n}} \big( \ell_{\mathbf{a}}(g_\mathbf{c}(0)) - \ell_{\mathbf{b}}(g_\mathbf{c}(0)) \big) + \mathcal{O}(\lambda^{-\mathfrak{n}}) .
\end{align*}
It follows that, for $t$ large enough, we can write, using the tree lemma:
\begin{align*}
& \lambda^{-2n-2\mathfrak{n}} \#\{ (\mathbf{a},\mathbf{b},\mathbf{c}) \in \mathcal{A}^{2n+2\mathfrak{n}} \ | \ \exists \theta, |\psi_{\mathbf{a},\mathbf{b},\mathbf{c}}^{(2)}(\theta)| \leq 1\}\\
&\quad \leq \lambda^{-2n-2\mathfrak{n}} \#\{ (\mathbf{a},\mathbf{b},\mathbf{c}) \in \mathcal{A}^{2n+2\mathfrak{n}} \ , \ \lambda^{\mathfrak{n}} \big| \ell_{\mathbf{a}}(g_\mathbf{c}(0)) - \ell_{\mathbf{b}}(g_\mathbf{c}(0)) \big| \leq 2\} \\
&\quad \leq \lambda^{-n-2\mathfrak{n}} \sum_{(\mathbf{b},\mathbf{c}) \in \mathcal{A}^{n+2\mathfrak{n}}} \lambda^{-n} \#\{ \mathbf{a} \in \mathcal{A}^n, \ \ell_\mathbf{a}(g_\mathbf{c}(0)) \in [\ell_\mathbf{b}(g_\mathbf{c}(0))-2\lambda^{-\mathfrak{n}},\ell_\mathbf{b}(g_\mathbf{c}(0))+2\lambda^{-\mathfrak{n}} ] \} \\
&\quad \leq 2 \lambda \Big( \lambda^{- \gamma_{sep} \mathfrak{n}/2} + \lambda^{-n \gamma_{sep}} \Big) \lesssim t^{-\gamma_{sep}/40}.
\end{align*}
Since also the van der Corput contribution is $\lesssim t^{-1/16}$, we obtain
$$
\int_{\mathbb S^1}
\left|\mathcal L_{\lambda,i\Phi(t,m)}^{\langle n_\beta(t)\rangle}h\right|^2\,d\theta
\lesssim
t^{-\rho}\|h\|_{1,t}^2,
\qquad
\rho:=\min\left\{\frac1{16},\frac{\gamma_{sep}}{40}\right\}>0.
$$
Thus $\rho$ depends only on $\varphi$.
\end{proof}

\end{document}
